% ================================================================
\section{Listas de prioridades e heaps binários}\label{sec:heap}
% ================================================================

\subsection{Lista de prioridades}
Tabela em que cada dado tem uma \textbf{prioridade}; operações: (1) \emph{seleção} do de maior prioridade; (2) \emph{inserção}; (3) \emph{remoção} do de maior prioridade; (4) \emph{alteração} de prioridade; (5) \emph{construção} a partir de um conjunto inicial. (Fusão entra com os heaps avançados.)

\begin{center}
\small
\begin{tabular}{@{}lccc@{}}
\toprule
Operação & Lista \textbf{não} ordenada & Lista ordenada & Heap\\
\midrule
seleção & $\Theta(n)$ & $\Theta(1)$ & $O(1)$\\
inserção & $\Theta(1)$ & $\Theta(n)$ & $O(\log n)$\\
remoção & $\Theta(n)$ & $\Theta(1)$ & $O(\log n)$\\
alteração & $\Theta(n)$ & $\Theta(n)$ & $O(\log n)$\\
construção & $\Theta(n)$ & $\Theta(n\log n)$ & $O(n)$\\
\bottomrule
\end{tabular}
\end{center}
As duas listas têm alguma operação de atualização $\Omega(n)$; o heap equilibra tudo em $O(\log n)$.

\subsection{Definição e representação}
\begin{defi}[Heap -- convenção do curso]
Um \textbf{heap} (de máximo) é uma lista linear $s_1,\dots,s_n$ tal que
\[
s_i\le s_{\floor{i/2}}\qquad\text{para todo }1<i\le n .
\]
(Heap de \emph{mínimo}: $s_i\ge s_{\floor{i/2}}$.)
\end{defi}
Visualiza-se como uma \textbf{árvore binária completa} numerada por níveis, da esquerda para a direita: o nó $i$ tem pai $\floor{i/2}$, filho esquerdo $2i$ e direito $2i+1$ (se $\le n$). A condição diz: \emph{cada nó tem prioridade $\ge$ à dos filhos}. Não há ponteiros: é só um vetor.

\begin{center}
\begin{forest} bst
[95,label={[font=\tiny]above:1}
  [60,label={[font=\tiny]left:2} [39,label={[font=\tiny]left:4} [33,label={[font=\tiny]left:8}] [,vazio]] [28,label={[font=\tiny]right:5}]]
  [78,label={[font=\tiny]right:3} [66,label={[font=\tiny]left:6}] [70,label={[font=\tiny]right:7}]]]
\end{forest}
\qquad
\raisebox{3em}{\begin{tabular}{c|c|c|c|c|c|c|c}
1&2&3&4&5&6&7&8\\\hline 95&60&78&39&28&66&70&33
\end{tabular}}
\end{center}

\paragraph{Fatos de índice} (usados em todas as provas):
\begin{itemize}
  \item o nó $i$ está no nível $\floor{\log_2 i}+1$ (raiz no nível 1); a árvore tem $\floor{\log_2 n}+1$ níveis;
  \item as \textbf{folhas} são exatamente as posições $\floor{n/2}+1,\dots,n$ (pois $i$ tem filho sse $2i\le n$); há $\ceil{n/2}$ folhas;
  \item convenção de altura do Szwarcfiter/slides: \textbf{folhas têm altura 1}.
\end{itemize}

\subsection{Subir e descer}
\begin{minipage}[t]{0.47\linewidth}
\begin{pseudo}{subir$(i)$ -- após \emph{aumentar} a prioridade de $i$}
\begin{algorithmic}
\Procedure{subir}{$i$}
  \State $j := \floor{i/2}$
  \If{$j\ge1$}
    \If{$T[i].chave > T[j].chave$}
      \State $T[i]\Leftrightarrow T[j]$
      \State \Call{subir}{$j$}
    \EndIf
  \EndIf
\EndProcedure
\end{algorithmic}
\end{pseudo}
\end{minipage}\hfill
\begin{minipage}[t]{0.5\linewidth}
\begin{pseudo}{descer$(i,n)$ -- após \emph{diminuir} a prioridade de $i$}
\begin{algorithmic}
\Procedure{descer}{$i,n$}
  \State $j := 2i$
  \If{$j\le n$}
    \If{$j<n$ \textbf{e} $T[j+1].chave > T[j].chave$}
      \State $j := j+1$ \Comment{maior filho}
    \EndIf
    \If{$T[i].chave < T[j].chave$}
      \State $T[i]\Leftrightarrow T[j]$
      \State \Call{descer}{$j,n$}
    \EndIf
  \EndIf
\EndProcedure
\end{algorithmic}
\end{pseudo}
\end{minipage}

\medskip
Ambos percorrem um caminho da árvore: $O(\log n)$. \textbf{Descer escolhe o maior filho} (se trocasse com o menor, o maior filho ficaria abaixo de um valor menor).

\subsection{Inserção, remoção, alteração}
\begin{pseudo}{Inserção, remoção do máximo e alteração de prioridade}
\begin{algorithmic}
\Procedure{inserir}{$novo$} \Comment{$M$ = capacidade}
  \If{$n<M$} \State $T[n+1]:=novo$;\ $n:=n+1$;\ \Call{subir}{$n$}
  \Else\ \textbf{overflow}
  \EndIf
\EndProcedure
\Procedure{remover}{}
  \If{$n\ne0$} \State agir$(T[1])$;\ $T[1]:=T[n]$;\ $n:=n-1$;\ \Call{descer}{$1,n$}
  \Else\ \textbf{underflow}
  \EndIf
\EndProcedure
\Procedure{alterar}{$i, p$}
  \State $antiga := T[i].chave$;\ $T[i].chave := p$
  \If{$p > antiga$} \State \Call{subir}{$i$} \Comment{aumentou $\Rightarrow$ sobe}
  \Else \State \Call{descer}{$i,n$} \Comment{diminuiu $\Rightarrow$ desce}
  \EndIf
\EndProcedure
\end{algorithmic}
\end{pseudo}

\begin{cuidado}[Erro clássico (monitoria de revisão, Q5)]
Um aluno escreveu: ``$H[i]:=p$; \textbf{se} $H[i]>H[\floor{i/2}]$ \textbf{então} Descer$(i)$ \textbf{senão} Subir$(i)$''. Está \textbf{invertido}: se o novo valor ficou maior que o pai, ele viola a propriedade \emph{com o pai} e precisa \textbf{subir}; se não ficou maior que o pai, a única violação possível é com os filhos, e ele precisa \textbf{descer}. Versão correta: \textbf{se} $H[i]>H[\floor{i/2}]$ (com $i>1$) \textbf{então} Subir$(i)$ \textbf{senão} Descer$(i,n)$.
\end{cuidado}

\begin{exemplo}[{Traços conferidos -- heap $H=[95,60,78,39,28,66,70,33]$}]\leavevmode
\begin{itemize}
  \item \textbf{Aumentar} o nó 6 de 66 para 98: sobe trocando com 78 (pos.\ 3) e com 95 (pos.\ 1): $[98,60,95,39,$\allowbreak$28,78,70,33]$.
  \item \textbf{Remover a raiz}: coloca-se o último (33) na raiz, $[33,60,78,39,28,66,70]$; desce trocando com 78 (maior filho) e depois com 70: $[78,60,70,39,28,66,33]$.
  \item \textbf{Diminuir} a raiz de 95 para 37: desce trocando com 78 e com 70: $[78,60,70,39,$\allowbreak$28,66,37,33]$.
  \item \textbf{Inserir} 73: entra na pos.\ 9 (filho de 39), sobe trocando com 39 e com 60: $[95,73,78,60,$\allowbreak$28,66,70,33,39]$.
\end{itemize}
\end{exemplo}

\subsection{Construção em tempo linear}
\textbf{1ª tentativa} (inserções sucessivas): \textbf{para} $i:=2,\dots,n$ \textbf{faça} subir$(i)$ --- custa $O(n\log n)$ (e $\Theta(n\log n)$ no pior caso, p.ex.\ entrada crescente).

\textbf{2ª tentativa} (\emph{arranjar}): as folhas ($i>\floor{n/2}$) já são heaps; corrige-se de baixo para cima.
\begin{pseudo}{arranjar$(n)$}
\begin{algorithmic}
\Procedure{arranjar}{$n$}
  \For{$i := \floor{n/2}$ \textbf{decr.\ até} $1$} \State \Call{descer}{$i,n$} \EndFor
\EndProcedure
\end{algorithmic}
\end{pseudo}
\emph{Correção:} invariante ``antes de processar $i$, as subárvores de raízes $i+1,\dots,n$ são heaps''. Ao chamar descer$(i,n)$, as duas subárvores dos filhos de $i$ já são heaps, e descer transforma a subárvore de $i$ em heap.

\begin{lema}[Lista de heaps, ex.\ 6]\label{lema:alturas}
Numa árvore binária completa com $n$ nós (folhas com altura 1), no máximo $\ceil{n/2^j}$ nós têm altura $j$.
\end{lema}
\begin{proof}
Numa árvore completa o caminho mais longo a partir de $i$ é o da esquerda: $i\to 2i\to 4i\to\cdots$. Logo o nó $i$ tem altura $\ge j$ sse seu descendente $j-1$ níveis abaixo existe, isto é, sse $i\,2^{j-1}\le n$, ou seja, $i\le\floor{n/2^{j-1}}$. Assim, com $x=n/2^{j-1}$ e $m=\floor{x}$,
\[
\#\{\text{altura}=j\}=\floor{x}-\floor{x/2}=m-\floor{m/2}=\ceil{m/2}\le\ceil{x/2}=\ceil{n/2^j}.\qedhere
\]
\end{proof}

\begin{teo}[Lema 6.1 do Szwarcfiter / slide]
arranjar$(n)$ roda em $O(n)$.
\end{teo}
\begin{proof}
descer$(i,n)$ num nó de altura $j$ faz no máximo $j-1$ trocas. Pelo Lema \ref{lema:alturas}, o total é
\[
\sum_{j=2}^{h}(j-1)\ceil*{\frac{n}{2^j}}\ \le\ \sum_{j\ge2}(j-1)\Bigl(\frac{n}{2^j}+1\Bigr)
\ =\ n\sum_{j\ge2}\frac{j-1}{2^j}+O(\log^2 n)\ =\ n\cdot 1+O(\log^2n)=O(n),
\]
pois $\sum_{j\ge1}\frac{j-1}{2^j}=\sum_{j\ge1}\frac{j}{2^j}-\sum_{j\ge1}\frac1{2^j}=2-1=1$. (O slide faz a mesma conta separando $j$ par e ímpar e usando $\frac{j}{2^{j+1}}\le\frac{j-1}{2^j}$.)
\end{proof}
Intuição: metade dos nós (folhas) não custa nada, $1/4$ custa $\le1$, $1/8$ custa $\le2$, \ldots; só a raiz custa $\log n$.

\begin{exemplo}[Traços do \emph{arranjar} (conferidos)]\leavevmode
\begin{itemize}
  \item Lista de heaps, ex.\ 5: $[18,25,41,34,$\allowbreak$14,10,52,50,48]$ ($n=9$, processa $i=4,3,2,1$).\\
  $i=4$: $34\leftrightarrow50$; $i=3$: $41\leftrightarrow52$; $i=2$: $25\leftrightarrow50$, depois $25\leftrightarrow48$; $i=1$: $18\leftrightarrow52$, depois $18\leftrightarrow41$.\\
  Resultado: $[52,50,41,48,$\allowbreak$14,10,18,34,25]$.
  \item Monitoria (Q4): $[20,15,30,40,10,25,5]$ processa $i=3,2,1$: $i=3$ nada; $i=2$: $15\leftrightarrow40$; $i=1$: $20\leftrightarrow40$. Resultado $[40,20,30,15,10,25,5]$.
  \item Monitoria (Q13): $[7,2,15,9,1,20,13]$: $i=3$: $15\leftrightarrow20$; $i=2$: $2\leftrightarrow9$; $i=1$: $7\leftrightarrow20$, $7\leftrightarrow15$. Resultado $[20,9,15,2,1,7,13]$ (o mesmo que por inserções sucessivas neste caso).
\end{itemize}
\end{exemplo}

\subsection{Heapsort}
\begin{pseudo}{Heapsort (Szwarcfiter, Alg.\ 7.6)}
\begin{algorithmic}
\Procedure{heapsort}{$T,n$}
  \State \Call{arranjar}{$n$} \Comment{$O(n)$}
  \For{$m := n$ \textbf{decrescendo até} $2$}
    \State $T[1]\Leftrightarrow T[m]$ \Comment{o máximo vai para o fim}
    \State \Call{descer}{$1, m-1$} \Comment{$O(\log m)$}
  \EndFor
\EndProcedure
\end{algorithmic}
\end{pseudo}
Complexidade $O(n)+\sum_{m}O(\log m)=\Theta(n\log n)$ em \emph{todos} os casos (mesmo com a entrada ordenada). \textbf{In-place} (só troca no vetor) e \textbf{não estável} (a troca da raiz com o último embaralha iguais). É ótimo entre os algoritmos por comparação. Ex.: $[18,25,41,34,14,10,52,50,48]\to[10,14,18,25,34,41,48,50,52]$.

% ------------------------------------------------------------------
\subsection{Propriedades de heap e suas demonstrações}\label{sec:heapprops}
\begin{prova}[Q3 da AV1: ``saber demonstrar propriedades de heap'']
Quase toda prova de heap usa só três ferramentas: (i) a definição $s_i\le s_{\floor{i/2}}$ + \textbf{transitividade ao longo de um caminho}; (ii) os \textbf{fatos de índice} (pai, filhos, folhas, níveis); (iii) \textbf{indução} (na altura, no nível, ou no número de nós). Para ``provar ou dar contraexemplo'', teste primeiro exemplos pequenos ($n\le 8$).
\end{prova}

\begin{prop}[Ancestral domina]\label{prop:ancestral}
Se $j$ é ancestral de $i$ num heap, então $s_i\le s_j$. Em particular: todo caminho da raiz a uma folha é não crescente, a \textbf{raiz é o máximo} e \textbf{cada subárvore é um heap}.
\end{prop}
\begin{proof}
O caminho de $i$ até $j$ é $i,\floor{i/2},\floor{i/4},\dots,j$; aplicando a definição em cada aresta e a transitividade de $\le$, $s_i\le s_{\floor{i/2}}\le\cdots\le s_j$. A raiz é ancestral de todos. A condição do heap numa subárvore é a mesma condição restrita a ela.
\end{proof}

\begin{prop}[O mínimo está numa folha]
Num heap de máximo com chaves distintas, o menor elemento está numa folha, isto é, numa posição $>\floor{n/2}$.
\end{prop}
\begin{proof}
Se $i$ não é folha, tem filho $2i$ e $s_{2i}<s_i$ (distintas); logo $s_i$ não é o mínimo. (Com repetições: \emph{algum} mínimo está numa folha.)
\end{proof}

\begin{prop}[Onde está o $k$-ésimo maior]
Com chaves distintas, o $k$-ésimo maior elemento está num nível $\le k$, isto é, numa posição $\le 2^k-1$. Em particular o 2º maior é filho da raiz (posição 2 ou 3), e o 3º maior está nas posições $2,\dots,7$.
\end{prop}
\begin{proof}
Pela Prop.~\ref{prop:ancestral}, todos os ancestrais de $x$ são maiores que $x$. Se $x$ é o $k$-ésimo maior, só existem $k-1$ elementos maiores, então $x$ tem no máximo $k-1$ ancestrais, ou seja, está no nível $\le k$. As posições dos níveis $1..k$ são $1,\dots,2^k-1$.
\end{proof}

\begin{prop}[Altura]
Um heap com $n$ elementos tem altura $\floor{\log_2 n}+1$ (em níveis) $=\Theta(\log n)$; por isso subir/descer/inserir/remover são $O(\log n)$.
\end{prop}
\begin{proof}
O último nó, $n$, está no nível $\floor{\log_2n}+1$, e numa árvore completa todos os níveis anteriores estão cheios: $2^{h-1}\le n\le 2^h-1$, logo $h=\floor{\log_2n}+1$.
\end{proof}

\begin{prop}[Relação com ordenação]
Toda lista ordenada de forma não crescente é um heap; a recíproca é falsa (ex.: $[3,1,2]$). Por isso construir um heap é mais fácil que ordenar ($O(n)$ vs.\ $\Omega(n\log n)$) --- e, como heapsort ordena usando construção + $n$ remoções, não existe heap com construção \emph{e} remoção ambas $o(\log n)$ amortizado \emph{por comparação} (senão ordenaríamos em $o(n\log n)$).
\end{prop}

\begin{prop}[Troca de dois elementos -- Lista de heaps, ex.\ 2 e 3]
Seja $S$ um heap e $i<j$.
\begin{enumerate}[label=(\alph*)]
  \item Se $s_i<s_j$, trocar $s_i$ com $s_j$ \textbf{não} produz necessariamente um heap.
  \item Se $s_i>s_j$, trocar $s_i$ com $s_j$ \textbf{não} produz necessariamente um heap.
\end{enumerate}
\end{prop}
\begin{proof}[Contraexemplos (verificados)]
(a) $S=[10,9,1,8,7,0,0,6]$ é heap; $i=3$ ($1$), $j=4$ ($8$). Após a troca, $[10,9,8,1,$\allowbreak$7,0,0,6]$: a posição 8 ($6$) tem pai na posição 4 ($1$) e $6>1$. Intuição: o valor pequeno vai para a posição $j$, cujos filhos podem ser maiores que ele.
(b) $S=[10,9]$, $i=1$, $j=2$: após a troca $[9,10]$, filho maior que o pai. Intuição: o valor pequeno sobe para $i$ e pode ficar abaixo de um descendente maior.
\end{proof}

\begin{prop}[Número de heaps -- Lista de heaps, ex.\ 4]
Seja $H(n)$ o número de permutações de $1,\dots,n$ que são heaps (vale o mesmo para min-heap). Então $H(0)=H(1)=1$ e
\[
H(n)=\binom{n-1}{L}\,H(L)\,H(n-1-L),
\]
onde $L$ é o número de nós da subárvore esquerda da raiz: com $h=\floor{\log_2 n}$ e $u=n-(2^h-1)$ nós no último nível, $L=(2^{h-1}-1)+\min\{u,2^{h-1}\}$. Valores: $1,1,2,3,8,20,80,210,896,3360$ para $n=1..10$ (conferido por força bruta até $n=8$).
\end{prop}
\begin{proof}
A raiz é necessariamente $n$ (máximo). O formato da árvore é fixo, então $L$ é determinado por $n$. Escolhem-se quais $L$ dos $n-1$ valores restantes vão para a esquerda ($\binom{n-1}{L}$ modos); cada lado precisa ser um heap sobre seus valores, e o número de heaps só depende da quantidade de valores ($H(L)$ e $H(n-1-L)$).
\end{proof}

\begin{prop}[Correção de subir e descer]
(a) Se $T$ era heap e a chave de $i$ \emph{aumentou}, após subir$(i)$ $T$ é heap. (b) Se as subárvores dos filhos de $i$ são heaps (em particular se $T$ era heap e a chave de $i$ \emph{diminuiu}), após descer$(i,n)$ a subárvore de $i$ é heap.
\end{prop}
\begin{proof}[Esboço]
(a) Invariante: a única violação possível é entre o nó corrente e seu pai. Ao trocar com o pai $j$ (quando $s_i>s_j$), o antigo $s_j$ desce para $i$ e é $\ge$ aos filhos de $i$ (eram seus descendentes); o novo valor em $j$ é maior que o antigo, então não viola com os filhos de $j$. Repete-se para cima. (b) Invariante: a única violação é entre o nó corrente e seus filhos. Troca-se com o \emph{maior} filho: ele passa a dominar o outro filho e o valor que desceu. A violação desce um nível; termina numa folha ou quando não há violação.
\end{proof}

\begin{chave}
Heap = vetor + ``pai $\ge$ filho''. Subir após aumento, descer após diminuição (maior filho!). Inserir = pôr no fim + subir; remover = último na raiz + descer. Construção linear pelo \emph{arranjar} (de $\floor{n/2}$ até $1$) graças a $\le\ceil{n/2^j}$ nós de altura $j$. Heapsort $\Theta(n\log n)$, in-place, não estável.
\end{chave}
